% % % % % % % % % % % % % % % % % % % % % % % % % % % % % % % % % % % % % % % % % % % % % % % % % % % % % % % % % % % % % % % %
% SAT_STYLE_MATH_MCQS_FR.tex
% 1_04_fr.tex
%
% encoding: UTF-8
% EOL: LF
%
% format: LaTeX
% indent: spaces (2)
% width: 127
% % % % % % % % % % % % % % % % % % % % % % % % % % % % % % % % % % % % % % % % % % % % % % % % % % % % % % % % % % % % % % % %
La numération en base heptadécimale permet d'écrire les entiers avec les dix chiffres usuels et sept lettres supplémentaires %
\texttt{A}, \dots, \texttt{G}: \texttt{A} représente $10$, \texttt{B} représente $11$, \dots, \texttt{G} représente $16$, %
\texttt{10} représente $17$, \texttt{1G} représente $33$, et ainsi de suite. \\
\\
\noindent%
Nous nous intéressons ici à l'entier 
%
\begin{equation}
  1729 = 1^3 + 12^3 = 9^3 + 10^3,
\end{equation} 
%
connu sous le nom de \emph{nombre de Hardy--Ramanujan}, ou \emph{taxicab number}. C'est le plus petit entier pouvant s'exprimer 
comme somme $a^3 + b^3$\;\,($0 <a^3 <b^3$) de deux manières différentes. La question est: comment s'écrit%
%
\begin{equation*}
    \Biggl(\Bigl(\bigl(1729^{1729} \bmod{17}\bigr)^{1729} \bmod{17}\Bigr)^{1729} \bmod{17}\Biggr) \cdots
\end{equation*}
%
dans le système heptadécimal?
%
\begin{enumerate}
  \item $1$
  \item $8$
  \item $B$
  \item $C$
  \item $G$
\end{enumerate}
%
\begin{proof}
En notant $X$ ce nombre, nous avons d'après l'énoncé
%
\begin{equation}
    X^{1729} \equiv X \pmod{17}.
  \end{equation}
%
$X=1$ peut passer pour une réponse évidente mais ce n'est pas la bonne! Comme 
%
\begin{equation}
  1729 \equiv 12 \pmod{17}
\end{equation}
%
(car $1729=101\times 17+12$), on a %
%
\begin{equation}
  1729^{1729} \equiv 12^{1729} \pmod{17}.
\end{equation}
%
Nous démontrons de quatre façons différentes que 
%
\begin{equation}
  12^{1729}\equiv 12 \pmod{17}.
\end{equation}
%
et donc que 
%
\begin{align}
    \bigl(12^{1729}\bigr)^{1729} &\equiv 12^{1729} \equiv 12 \pmod{17}, \\
    \Bigl(\bigl(1729^{1729}\bigr)^{1729}\Bigr)^{1729} &\equiv 12^{1729} \equiv 12 \pmod{17},
    %&\hspace{0.25em} \vdots\nonumber
\end{align}
%
et ainsi de suite. Comme $0 \leq 12 <17$, $12$ est le reste recherché. Dans le système heptadécimal, il s'écrit $C$. Réponse %
\textbf{C}.\\
\\
Notre première preuve mobilise le Petit Théorème de Fermat, la deuxième son équivalent algorithmique qu'est l'exponentiation %
rapide. Sans l'un de ces outils, résoudre les congruences est beaucoup plus difficile et dépend de cas particuliers. Les deux %
autres preuves sont en [\ref{annex:1.29}]. Elles servent de prétexte à un passage en revue des autres concepts et méthodes %
importants en arithmétique modulaire. Par contraste, leur difficulté met en valeur la force du Petit Théorème de Fermat.
%
% % % % % % % % % % % % % % % % % % % % % % % % % % % % % % % % % % % % % % % % % % % % % % % % % % % % % % % % % % % % % % % %
% First variant (Fermat)
% % % % % % % % % % % % % % % % % % % % % % % % % % % % % % % % % % % % % % % % % % % % % % % % % % % % % % % % % % % % % % % %
\paragraph{Avec le petit théorème de Fermat} Si $p$ est premier, le petit théorème de Fermat implique %
%
\begin{equation}
  n^{p-1} \equiv 1 \pmod{p}
\end{equation}
%
pour tout entier $n \nequiv 0 \pmod{p}$. Donc, comme $17$ est premier, 
%
\begin{equation}
  12^{16} \equiv 1 \pmod{17}
\end{equation}
%
en application de ce théorème. Sachant que $1729 = q \times 16 + 1$ ($q=108$), nous obtenons 
%
\begin{align}
  12^{1729} &= (12^{16})^q \times 12^1 \\
  &\equiv 1^{q} \cdot 12 \pmod{17} \\
  &\equiv 12 \pmod{17}; 
\end{align}
%
ce qui termine cette première démonstration.
%
% % % % % % % % % % % % % % % % % % % % % % % % % % % % % % % % % % % % % % % % % % % % % % % % % % % % % % % % % % % % % % % %
% Second variant: rapid exp
% % % % % % % % % % % % % % % % % % % % % % % % % % % % % % % % % % % % % % % % % % % % % % % % % % % % % % % % % % % % % % % %
\paragraph{Par une exponentiation rapide} Dans ce qui suit, le carré modulo $17$ exprimé dans une ligne est repris par la %
ligne suivante:
%
\begin{align}
  12^2 &= 144 \equiv 8 \pmod{17}, \\
  12^4 &\equiv 8^2  = 64 \equiv 13 \pmod{17},  \\
  12^8 &\equiv 13^2 = 169 \equiv \minus 1 \pmod{17}, \\
  &\hspace{0.25em}\vdots\nonumber
\end{align}
%
Sachant que $1729 = (2\times q) \times 8 + 1$ ($q=108$), nous obtenons 
%
\begin{align}
  12^{1729} &= (12^8)^{2\times q}  \times 12^1 \\
  &= (-1)^{2\times q}  \times 12^1 \\
  &\equiv 12 \pmod{17}.
\end{align}
%
La deuxième démonstration est ainsi achevée. 
\end{proof}